\begin{proposition}[nonvanishing and bases]
\label{prop:bcpr-generators-are-bases}
\Verified\ Together with the unit $\mathbf1\in H^{0,0}_{BC}(X)$, the entries
in \cref{eq:bcpr-generators} give bases of all the nonzero Bott--Chern groups
in \cref{eq:bcpr-nonzero-dimensions}:
$\mathsf A,\mathsf B$ form the basis in bidegree $(1,1)$ and
$\mathsf D,\overline{\mathsf D}$ form the basis in bidegree $(2,2)$;
each other listed bidegree is one-dimensional and is generated by its
displayed class.
\end{proposition}

\begin{proof}
Because $X$ is compact and connected, a smooth function annihilated by both
$\partial$ and $\bar\partial$ is constant.  There is no
$\partial\bar\partial$ image in bidegree $(0,0)$, so
$\mathbf1:=[1]_{BC}$ generates
$H^{0,0}_{BC}(X)=\mathbb C\mathbf1$.

The Dolbeault image of $\mathsf A$ is
$d_1\mathfrak g\neq0$ by injectivity in
\cref{eq:bcpr-hodge-row-one}.  Thus $\partial G$ is a nonzero
Bott--Chern class.  The Dolbeault image of $\mathsf B$ is zero because
$a^\dagger=\bar\partial\bar G$, while
$\mathsf B=\overline{\mathsf A}\neq0$.  Hence
$\mathsf A,\mathsf B$ are independent, and
\cref{eq:bcpr-nonzero-dimensions} makes them a basis.

The classes $d_1\mathfrak g$ and $\mathfrak q$ are independent in
$H^{1,1}_{\bar\partial}(X)$: the first restricts to zero on every smooth
torus fibre, while the second restricts to the nonzero class
by \cref{eq:bcpr-q-restriction,eq:bcpr-q-square-integral}.  Exactness of
\cref{eq:bcpr-hodge-row-one} therefore gives
\begin{equation}
 d_1\mathfrak q=[\partial Q]_{\bar\partial}\neq0.
\label{eq:bcpr-c-nonzero}
\end{equation}
Thus $\mathsf C\neq0$, and conjugation gives
$\overline{\mathsf C}\neq0$.

Set
\begin{equation}
 v=[G\wedge Q]_{\bar\partial}
 \in H^{1,2}_{\bar\partial}(X).
\label{eq:bcpr-v-class}
\end{equation}
Its fibre restriction is the class represented by
$\Gamma\wedge Q_F$.  This class is nonzero, rather than merely represented
by a nonzero invariant form: wedging with the holomorphic one-form
$d\zeta_2$ and using \cref{eq:bcpr-q-gamma-product} gives
\begin{equation}
 \int_{F_b}Q_F\wedge\Gamma\wedge d\zeta_2
 =\frac1{4y\Delta}
   \int_{F_b}d\zeta_1\wedge d\zeta_2\wedge
                    d\bar\zeta_1\wedge d\bar\zeta_2
 =-1
\label{eq:bcpr-q-gamma-serre-pairing}
\end{equation}
for the complex orientation: the middle integral equals
$-4y\Delta>0$.  Thus the displayed constant top form has nonzero integral.
Dolbeault--Serre duality on $F_b$ therefore gives
$[\Gamma\wedge Q_F]_{\bar\partial}\ne0$.  Equivalently, if
$Q_F\wedge\Gamma=\bar\partial_F\Psi$ with $\Psi$ of bidegree $(1,1)$,
then, since $\bar\partial_Fd\zeta_2=0$,
\[
 Q_F\wedge\Gamma\wedge d\zeta_2
 =\bar\partial_F(\Psi\wedge d\zeta_2),
\]
and Stokes' theorem would make the integral in
\cref{eq:bcpr-q-gamma-serre-pairing} zero.  Here
$Q_F\wedge\Gamma$ has bidegree $(1,2)$,
$d\zeta_2$ has bidegree $(1,0)$, and
$\Psi\wedge d\zeta_2$ has bidegree $(2,1)$; its $\partial_F$ is zero in
complex dimension two, so the displayed $\bar\partial_F$-boundary is a
$d$-boundary.  Thus $v\neq0$, and it spans
$H^{1,2}_{\bar\partial}(X)$.  Injectivity in
\cref{eq:bcpr-hodge-row-two} gives
\begin{equation}
 d_1v=[\partial(G\wedge Q)]_{\bar\partial}\neq0.
\label{eq:bcpr-d-nonzero}
\end{equation}
Hence $\mathsf D\neq0$, and conjugation gives
$\overline{\mathsf D}\neq0$.  The latter has zero Dolbeault image because
its representative is $\bar\partial$-exact, so it is independent of
$\mathsf D$; \cref{eq:bcpr-nonzero-dimensions} gives the asserted
$(2,2)$ basis.

The class $[Q^2]_{\bar\partial}$ has nonzero fibre restriction by
\cref{eq:bcpr-q-square-integral}.  Since $v$ spans
$H^{1,2}_{\bar\partial}(X)$, the image of $d_1$ is spanned by $d_1v$.
On a smooth fibre,
\[
 i_b^*(d_1v)
 =d_1[\Gamma\wedge Q_F]_{\bar\partial}
 =[\partial_F(\Gamma\wedge Q_F)]_{\bar\partial}=0,
\]
because $\Gamma$ and $Q_F$ have constant fibre coefficients.  Thus every
class in $\operatorname{im}(d_1\colon H^{1,2}_{\bar\partial}\to
H^{2,2}_{\bar\partial})$ restricts to zero.
Therefore $[Q^2]_{\bar\partial}$ is independent of $d_1v$.  Exactness of
\cref{eq:bcpr-hodge-row-two} yields
\begin{equation}
 d_1[Q^2]_{\bar\partial}
 =[\partial(Q^2)]_{\bar\partial}\neq0.
\label{eq:bcpr-e-nonzero}
\end{equation}
Hence $\mathsf E\neq0$, and conjugation gives
$\overline{\mathsf E}\neq0$.  Finally, $\int_X\nu=1$ proves
$\mathsf T\neq0$.  The dimensions in
\cref{eq:bcpr-nonzero-dimensions} complete the proof.
\end{proof}
